% !TEX root = ../../main.tex
% ============================================================================
% Part II; structural introduction.
%
% The Characteristic Datum Part opens with a deficiency. Classical
% conformal field theory recognizes ONE shadow coordinate: the central
% charge $c$, emitted by the Virasoro operator product expansion as
% the quartic coefficient of $T(z)T(w)$. Every textbook narrative of
% the rational landscape inherits this one-coordinate reading; every
% "structural" theorem ends at $c$.
%
% The structural formulation says the
% one-coordinate reading is a shadow. The Koszul depth $\kappa(\cA)$
% is the $r=2$ rung of an infinite shadow ladder
% $S_\bullet(\cA) = \{S_2, S_3, S_4, S_5, S_6, S_7, S_8, \ldots\}$;
% the pole order $p_{\max}(\cA)$, the chiral Euler character
% $\chi_{\VOA}(\cA)$, the strong-generator count $n_{\mathrm{strong}}(\cA)$,
% and the coset type $\mathrm{coset}(\cA)$ are four further invariants
% genuinely independent of $\kappa$. Their assembly
% $\varphi(\cA) = (p_{\max}, r_{\max}, \chi_{\VOA}, n_{\mathrm{strong}},
% \mathrm{coset})$ is the conditional reconstruction datum for the
% bar-complex finite-window structure of $\cA$
% (Theorem~\ref{thm:part-ii-fingerprint-reconstruction-surface});
% the familiar four-class assignment G/L/C/M is its lossy coarse
% projection onto the second coordinate.
%
% This chapter is the structural introduction to Part~II: a replacement of
% the single-coordinate framing to a fingerprint theorem schema with
% one structural invariant and seven shadows. The eight sections of the
% reconstituted Part~II are laid out, the five structural surfaces stated,
% and the dependency map between existing Part~II chapters and the
% structural theorems exhibited.
%
% Structure:
% \S0 Deficiency of the one-coordinate theory.
% \S1 The shadow tower (backbone).
% \S2 G/L/C/M as coarse projection (explicit witnesses).
% \S3 Infinite fingerprint classification.
% \S4 Motivic shadow tower.
% \S5 Higher $S_r$ coefficients through $r = 8$.
% \S6 Chiral moonshine as a fingerprint orbit.
% \S7 Higher-genus fingerprint.
% \S8 Structural endpoint: the completeness theorem.
%
% Reconstituted $2026$-$04$-$16$ (Raeez Lorgat).
% ============================================================================

\chapter{The shadow-tower fingerprint of a chiral algebra}
\label{ch:part-ii-platonic-introduction}

The Bar--Cobar Engine of Part~\ref{part:bar-cobar-engine} produces, for
each chart algebra $A_b = \End_\cC(b)$, a bar coalgebra $B(A_b)$.
Its ordered scalar residue tower is defined after a package
$\mathsf H_{\mathrm{res}}(A_b;X)$ specifies the residue complex,
collision maps, and scalar contractions.  The Virasoro operator
product expansion
\begin{equation}
 T(z)\,T(w) \;\sim\; \frac{c/2}{(z-w)^4}
 \;+\; \frac{2\,T(w)}{(z-w)^2}
 \;+\; \frac{\partial T(w)}{z-w}
\end{equation}
has central coefficient $c/2=\kappa(\Vir_c)$.  This coefficient is the
arity-two local input for the normalized scalar coordinate supplied by
$\mathsf H_{\mathrm{res}}$.  Higher coordinates require the transferred
residue construction.

The one-coordinate reading is a shadow. Every chiral algebra in
the standard landscape carries, in addition to $\kappa$, four
further structural invariants genuinely independent of $\kappa$ and
of each other:
\begin{equation}
 \varphi(\cA) \;=\; \bigl(\,p_{\max}(\cA),\; r_{\max}(\cA),\;
 \chi_{\VOA}(\cA),\; n_{\mathrm{strong}}(\cA),\;
 \mathrm{coset}(\cA)\,\bigr).
 \label{eq:part-ii-fingerprint}
\end{equation}
The first slot is the operator-product pole order; the second is
the depth of a constructed scalar residue tower; the third is the
Hodge-filtered chiral Euler character;
the fourth is the minimal strong-generator count of the
strong-generating set compatible with the PBW filtration; the fifth
is the isomorphism class of the commutant pair
$(\Com(H_\kappa, \cA), \cA/\Com)$. The five-tuple $\varphi(\cA)$
is the conditional reconstruction datum for the bar-complex
finite-window structure of $\cA$
(Theorem~\ref{thm:part-ii-fingerprint-reconstruction-surface}); the
familiar four-class labelling of the standard landscape as
$\mathrm{G}/\mathrm{L}/\mathrm{C}/\mathrm{M}$ is the lossy coarse
projection of $\varphi$ onto the $r_{\max}$ coordinate restricted
to $\kappa \ne 0$ (Theorem~\ref{thm:part-ii-quadrichotomy-coarse-projection}).

The exact Virasoro local calculation fixes the pole-order and
strong-generator slots.  Its typed fingerprint therefore begins
\begin{equation}
 \varphi_{\mathrm{loc}}(\Vir_c) \;=\;
 \bigl(\,4,\; r_{\max}^{\mathrm{res}},\; \chi_{\VOA}(\Vir_c),\; 1,\;
 \mathrm{coset}_{\mathrm{Vir}}\,\bigr),
\end{equation}
where $r_{\max}^{\mathrm{res}}$ is determined by a compatible
all-arity residue package.  The exact level-four vacuum Gram
determinant is
\begin{equation}
 \det G^{\Vir}_{4}(c) \;=\; \tfrac{1}{2}\,c^2(5c + 22).
\end{equation}
The multiplicity of this factor in higher scalar coordinates belongs
to the denominator package of
Chapter~\ref{ch:shadow-tower-higher-coefficients}.  The weight-$1/2$
McKay--Thompson
series of $V^\natural$ lift to universal elements of the twining
fingerprint
(Theorem~\ref{thm:part-ii-moonshine-via-fingerprint}); and the
structural endpoint
(Theorem~\ref{thm:part-ii-platonic-completeness}) assembles the
five slots into a single canonical moduli-theoretic object.

\medskip

Part~II reconstituted carries eight sections indexed by shadow
coordinate. The characteristic-datum chapters of the Bar--Cobar Engine
are retained; this structural introduction orients the reader by
laying out the deficiency, the ladder, the section architecture,
and the five theorems.

\section{Deficiency of the one-coordinate theory}
\label{sec:part-ii-deficiency}

Fix a chirally Koszul algebra $\cA$ in the standard landscape.
The classical extraction produces $c(\cA)$ (if $\cA$ has a
Virasoro sub-VOA) or $\kappa(\cA)$ (for every chirally Koszul
$\cA$ in the standard landscape with $\kappa \ne 0$, via the
bar-complex shadow $S_2$; the $\kappa = 0$ Feigin--Frenkel
companion is treated separately). Four independent obstructions show
this is not the full story.

\begin{enumerate}
 \item \textbf{Two distinct $\kappa$-isospectral chiral algebras.}
 Write $H_{\ell}$ for the Heisenberg algebra at level $\ell$ with
 $\kappa(H_{\ell}) = \ell$, and $V_{k}(\fsl_2)$ for the affine
 Kac--Moody algebra at level $k$ with
 $\kappa(V_{k}(\fsl_2)) = 3(k+2)/4$. The tuning
 $k = \tfrac{4}{3}\ell - 2$ produces two non-isomorphic algebras
 with identical $\kappa = \ell$. The two pole orders coincide at
 $p_{\max}(H_{\ell}) = p_{\max}(V_{k}(\fsl_2)) = 2$, so $p_{\max}$
 does \emph{not} separate the pair; the $r_{\max}$ slot does:
 $r_{\max}(H_{\ell}) = 2$ places $H_{\ell}$ in class $\mathrm{G}$,
 while $r_{\max}(V_{k}(\fsl_2)) = 3$ places $V_{k}(\fsl_2)$ in
 class $\mathrm{L}$. The Koszul depth alone fails to separate;
 the shadow depth is the first distinguishing coordinate.
 \item \textbf{Shadow depth $r_{\max}$ genuinely independent of
 $p_{\max}$.} Theorem~\ref{thm:pole-depth-independence} of
 Chapter~\ref{ch:infinite-fingerprint} exhibits explicit
 witnesses: $\beta\gamma$ has $(p_{\max}, r_{\max}) = (1, 4)$;
 $H_k$ has $(2, 2)$; $V_k(\fsl_2)$ has $(2, 3)$; $\Vir_c$ has
 $(4, \infty)$. The two invariants are independent; neither
 determines the other, and neither alone determines the class.
 \item \textbf{The motivic Galois group acts on the shadow
 tower.} The shadow coefficients
 $S_r(\Vir_c) \in \mathbb{Q}(c)$ at $r = 4, 5, 6, 7, 8$ are
 $\mathbb{Q}$-linear combinations of weight-$(r-2)$ multiple
 zeta values whose specializations at period maps are the
 classical rationals (Theorem~\ref{thm:shadow-tower-motivic-lift}).
 No classical reading of the central charge $c$ alone captures
 this motivic structure; the Tannakian lift requires the full
 tower $S_\bullet$, not just $S_2 = \kappa$.
 \item \textbf{Chiral moonshine fingerprints exceed the
 central-charge data.} The Monster vertex operator algebra
 $V^\natural$ has $c = 24$, yet the McKay--Thompson series
 $T_g(\tau)$ for $g \in M$ carry class-function data
 parametrized by the $194$ conjugacy classes of $M$; the
 fingerprint $\varphi(V^\natural)$ encodes, via the coset slot,
 the $\Lambda_{\mathrm{Leech}}^+$ orbifold structure that is
 invisible to $c$ alone
 (Theorem~\ref{thm:part-ii-moonshine-via-fingerprint}).
\end{enumerate}

The classical one-coordinate extraction recovers
$\kappa(\cA) = c(\cA)/2$ (for $\cA$ with Virasoro sub-VOA) or
$\kappa(\cA)$ (for every chirally Koszul $\cA$ in the standard
landscape with $\kappa \ne 0$); this is the $r = 2$ rung of a
shadow tower of depth $r_{\max} \in \{2, 3, 4, \infty\}$, and the
lossy projection of a fingerprint of five coordinates onto its
second slot. The
characteristic datum of Part~II is the fingerprint; the central
charge is its weight-$2$ shadow.

\section{The shadow tower}
\label{sec:part-ii-shadow-tower}

The structural backbone of Part~II is the ordered scalar residue
tower.  Its local inputs and construction obligations are recorded in
Chapter~\ref{ch:shadow-tower-higher-coefficients}.

\begin{definitionref}[Shadow obstruction tower;
 recalling Definition~\ref{def:shadow-obstruction-tower}]
\label{def:part-ii-shadow-tower-recall}
Let $\cA$ be a chirally Koszul chiral algebra with bar complex
$B(\cA) = T^c(s^{-1}\bar\cA)$.  Assume an ordered-residue package
$\mathsf H_{\mathrm{res}}(\cA;X)$ and let
$s_\cA$ be the transferred scalar Maurer--Cartan element.  Write
\begin{equation}
 s_\cA \;=\; \sum_{r\ge2}
 S_r(\cA;\mathsf H_{\mathrm{res}})e_r.
\end{equation}
The residue depth is
\begin{equation}
 r_{\max}^{\mathrm{res}}(\cA;\mathsf H_{\mathrm{res}})
 :=\sup\bigl\{r\ge2:S_r(\cA;\mathsf H_{\mathrm{res}})\ne0\bigr\}
 \;\in\; \{2, 3, 4, \ldots, \infty\}.
\end{equation}
\end{definitionref}

For Virasoro, the exact inputs are the operator product, the Ward
correlators, and the level-four Gram form.
Problems~\ref{thm:s6-virasoro-closed-form},~\ref{thm:s7-virasoro-closed-form},
and~\ref{thm:s8-virasoro-closed-form} ask for the corresponding
chain-level residue computations.

\section{G/L/C/M as coarse projection}
\label{sec:part-ii-glcm}

The four-class assignment $\mathrm{G}/\mathrm{L}/\mathrm{C}/\mathrm{M}$
on the finite standard-landscape atlas is indexed by shadow depth.
Under the fingerprint reading, this assignment is a lossy coarse
projection of the five-slot fingerprint, not a classification of
arbitrary chiral algebras.

\begin{theorem}[Quadrichotomy is coarse projection of
 fingerprint; \ClaimStatusConditional]
\label{thm:part-ii-quadrichotomy-coarse-projection}
Let $\cA$ be a chirally Koszul chiral algebra with $\kappa(\cA)
\ne 0$. The four-class assignment
$\mathrm{class}(\cA) \in \{\mathrm{G}, \mathrm{L}, \mathrm{C},
\mathrm{M}\}$ is the lossy coarse projection
\begin{equation}
 \mathrm{class}(\cA) \;=\;
 \Pi_{\mathrm{coarse}}\bigl(\varphi(\cA)\bigr)
 \;=\;
 \begin{cases}
 \mathrm{G} & r_{\max}(\cA) = 2, \\
 \mathrm{L} & r_{\max}(\cA) = 3, \\
 \mathrm{C} & r_{\max}(\cA) = 4, \\
 \mathrm{M} & r_{\max}(\cA) = \infty,
 \end{cases}
\end{equation}
with explicit witnesses as follows.
\begin{itemize}
 \item G: Heisenberg $H_k$ with $\varphi(H_k) = (2, 2,
 \chi^H, 1, \mathrm{coset}_H)$.
 \item L: affine Kac--Moody $V_k(\fsl_2)$ with
 $\varphi(V_k(\fsl_2)) = (2, 3, \chi^{KM}, 3, \mathrm{coset}_{KM})$.
 \item C: $\beta\gamma$-system with $\varphi(\beta\gamma) =
 (1, 4, \chi^{\beta\gamma}, 2, \mathrm{coset}_{\beta\gamma})$.
 \item M: the Virasoro row with
 $\varphi(\Vir_c)=(4,r_{\max}^{\mathrm{res}},
 \chi^{\Vir}_c,1,\mathrm{coset}_{\Vir})$; its placement at
 $r_{\max}^{\mathrm{res}}=\infty$ is the all-arity residue
 construction problem.
\end{itemize}
At $\kappa(\cA) = 0$, a fifth stratum $\mathrm{FF}$
\textup{(}Feigin--Frenkel\textup{)} appears as a fully canonical
companion, realized on the critical-level affine Kac--Moody
$V_{-h^\vee}(\fg)$ with
$\varphi(V_{-h^\vee}(\fg)) = (2, \star, \chi^{FF}, 3,
\mathrm{coset}_{FF})$, where the $r_{\max}$ slot is replaced by
the Feigin--Frenkel centre stratification datum. The five-class
stratification $\mathrm{G}/\mathrm{L}/\mathrm{C}/\mathrm{M}/\mathrm{FF}$
is the full coarse projection of $\varphi$.
\end{theorem}

\begin{proof}
The fingerprint $\varphi$ records the finite-window row data by
Theorem~\ref{thm:fingerprint-completeness}; the conditional
reconstruction surface is
Theorem~\ref{thm:part-ii-fingerprint-reconstruction-surface}. Class
assignment via shadow depth is the $r_{\max}$-separation theorem
\textup{(}Theorem~\ref{thm:pole-depth-independence}\textup{)}
together with the intrinsic shadow-tower reading of
Theorem~\ref{thm:fingerprint-reconstruction-equivalence}. The explicit witnesses are
computed in Appendix~\ref{app:nonlinear-modular-shadows}
\textup{(}Heisenberg, $V_k(\fsl_2)$\textup{)},
Chapter~\ref{chap:beta-gamma} \textup{(}$\beta\gamma$\textup{)}, and
Chapter~\ref{chap:w-algebras} \textup{(}$\Vir_c$\textup{)}; the
$r_{\max}$ coordinate separates the four cases by Chapter
Theorem~\ref{thm:pole-depth-independence}. The Feigin--Frenkel
companion stratum $\mathrm{FF}$ is treated in
Chapter~\ref{chap:kac-moody} at the critical-level discussion; the
$r_{\max}$ slot is replaced by the opers stratification of
Feigin--Frenkel, and the remaining slots descend unchanged.
\end{proof}

The quadrichotomy advertisement is strictly weaker than the
fingerprint it projects: two algebras with identical class can
differ in $(p_{\max}, \chi_{\VOA}, n_{\mathrm{strong}},
\mathrm{coset})$; examples are in Chapter
Section~\ref{sec:fingerprint-non-injective-on-class}.

\section{Infinite fingerprint classification}
\label{sec:part-ii-fingerprint}

The five-slot fingerprint
$\varphi(\cA) \in \mathcal{F}$ is a finite-window datum. It becomes
a bar-complex reconstruction invariant only on the residual-free
fingerprint-reconstruction surface.

\begin{theorem}[Conditional fingerprint reconstruction surface;
 \ClaimStatusConditional]
\label{thm:part-ii-fingerprint-reconstruction-surface}
\textup{[}Type signature:
$(\mathrm{Open},\mathrm{fingerprint/bar\ reconstruction},
\mathrm{level}=1\leftrightarrow2,
\mathrm{hypothesis\ package}=\mathrm{standard\ landscape}
+\mathrm{slot\ independence}
+\mathrm{bar\text{-}intrinsic\ MC}
+\mathrm{residual\text{-}free\ finite\text{-}window\ reconstruction})$.\textup{]}
For any two chirally Koszul chiral algebras $\cA, \cB$ in the
standard landscape, an isomorphism
$B(\cA) \simeq B(\cB)$ of augmented coaugmented conilpotent
$E_1$-chiral coassociative coalgebras implies
$\varphi(\cA) = \varphi(\cB)$. Conversely, if the
residual-free fingerprint-reconstruction package holds for the pair
$(\cA,\cB)$---i.e.\ the finite-row data, higher braces, ordered
coinvariants, and OPE coefficients forgotten by the projection
$\Pi_{\mathrm{fp}}$ are already recovered from the five slots in the
chosen completed landscape row---then
$\varphi(\cA)=\varphi(\cB)$ implies
$B(\cA) \simeq B(\cB)$. In particular, the five coordinates are
independent
\textup{(}Theorem~\ref{thm:pole-depth-independence},
Proposition~\ref{prop:chi-strong-coset-independence}\textup{)} and
jointly reconstructive only under the stated residual-free package.
The projection
$\Pi_{\mathrm{coarse}}: \mathcal{F} \to \{\mathrm{G}, \mathrm{L},
\mathrm{C}, \mathrm{M}, \mathrm{FF}\}$ is a genuine quotient with
non-trivial fibres: e.g.\ all $V_k(\fsl_n)$ at generic $k$ for
$n \ge 2$ share class $\mathrm{L}$ but carry distinct
fingerprints via $n_{\mathrm{strong}}(V_k(\fsl_n)) = \dim(\fsl_n)
= n^2 - 1$.
\end{theorem}

\begin{proof}
Forward \textup{(}isomorphism implies equality of fingerprint\textup{)}:
each slot is a bar-complex-isomorphism invariant. $p_{\max}$ is
the maximum pole order visible in the collision differential of
$B(\cA)$; $r_{\max}$ is the shadow-tower stabilization index; the
Hodge-filtered chiral Euler character $\chi_{\VOA}(\cA)$ is a
bigraded supertrace on $B(\cA)$; $n_{\mathrm{strong}}$ is the
minimal strong-generator count compatible with the PBW
filtration on $B(\cA)$; $\mathrm{coset}$ is the iso class of the
commutant pair realized inside $B(\cA)$.

Reverse \textup{(}equality of fingerprint implies bar-complex
isomorphism on the residual-free surface\textup{)}:
the hypothesis package of
Theorem~\ref{thm:fingerprint-reconstruction-equivalence} supplies a
reconstruction functor
\begin{equation}
 \Phi_{\mathrm{recon}}\colon \mathcal{F}
 \;\longrightarrow\; E_1\text{-coAlg}^{\mathrm{chir},
 \mathrm{aug}, \mathrm{conil}}_{/X}
\end{equation}
such that $\Phi_{\mathrm{recon}}(\varphi(\cA)) \simeq B(\cA)$ for
all $\cA$ in the stated residual-free completed row; the five slots determine,
in order, the collision-differential structure, the shadow
stabilization pattern, the Hodge bigrading, the generator-relation
data, and the commutant decomposition, with successive slots
determining the data not fixed by earlier slots. Independence
of slots is
Theorem~\ref{thm:pole-depth-independence} and
Proposition~\ref{prop:chi-strong-coset-independence}.
\end{proof}

The fingerprint theorem is the structural replacement of the
four-class narrative: one conditional reconstruction surface, one
coarse projection, one canonical stratification.

\section{Motivic residue comparison}
\label{sec:part-ii-motivic}

The exact Virasoro Ward and Gram calculations provide the local
source for a motivic comparison.  The ordered-residue package
constructs the scalar class, and the motivic package of
Definition~\ref{def:vmp-motivic-lift} maps its chain representative
to a motivic period complex.

\begin{problem}[Construct the motivic Virasoro residues through arity \(8\)]
\ClaimStatusOpen
\label{thm:part-ii-motivic-purity-virasoro-r-le-8}
\emph{Type signature:}
\textup{(Open/motivic quadrants, ordered residue and motivic-period
presentations, levels \(2\leftrightarrow5\),
\(\mathsf H_{\mathrm{res}}^{\le8}
+\mathsf H_{\mathrm{mot}}^{\le8}\)).}

Construct compatible packages
\(\mathsf H_{\mathrm{res}}^{\le8}(\Vir_c;X)\) and
\(\mathsf H_{\mathrm{mot}}^{\le8}\).  For every
\(2\le r\le8\), produce a class
\[
 S_r^{\mot}(\Vir_c)
 \in\mathcal P_r^{\mot}\otimes\mathbb Q(c)
\]
whose period is the ordered scalar residue.  Determine its reduced
coaction, its Tate twists, and the divisor of its rational-function
specialization,
\begin{equation}
 D_r^{\mathrm{res}}(c)
 :=\operatorname{div}_c
 \bigl(S_r(\Vir_c;\mathsf H_{\mathrm{res}})\bigr).
 \label{eq:part-ii-motivic-denominator}
\end{equation}
Tate factorization is governed by
\(\mathsf H_{\mathrm{Tate}}^{\le8}\); denominator control is governed
by \(\mathsf H_{\mathrm{den}}^{\le8}\).
\end{problem}

\section{Higher residue coordinates}
\label{sec:part-ii-higher-s-r}

\begin{problem}[Compute the Virasoro scalar coordinates through arity \(8\)]
\ClaimStatusOpen
\label{prop:part-ii-virasoro-shadow-through-r-8}
\emph{Type signature:}
\textup{(Open quadrant, completed ordered residue-bar presentation,
level \(2\), \(\mathsf H_{\mathrm{res}}^{\le8}\)).}
Assume \(\mathsf H_{\mathrm{res}}^{\le8}(\Vir_c;X)\) and write
\(s_{\Vir_c}\) for its transferred scalar Maurer--Cartan element.
The coordinates to be computed are
\begin{align}
 S_6(\Vir_c;\mathsf H_{\mathrm{res}})
 &=\operatorname{pr}_6(s_{\Vir_c}),
 \label{eq:s6-virasoro}\\
 S_7(\Vir_c;\mathsf H_{\mathrm{res}})
 &=\operatorname{pr}_7(s_{\Vir_c}),
 \label{eq:s7-virasoro}\\
 S_8(\Vir_c;\mathsf H_{\mathrm{res}})
 &=\operatorname{pr}_8(s_{\Vir_c}).
 \label{eq:s8-virasoro}
\end{align}
Derive these coordinates from collision-compatible chain
representatives and verify them through a second contraction model.
Problems~\ref{thm:s6-virasoro-closed-form},~\ref{thm:s7-virasoro-closed-form},
and~\ref{thm:s8-virasoro-closed-form} record the corresponding
construction obligations.
\end{problem}

\section{Chiral moonshine via fingerprint}
\label{sec:part-ii-moonshine}

The moonshine correspondence $V^\natural \leftrightarrow
\mathrm{Monster}$ lifts to a fingerprint-twining correspondence.

\begin{theorem}[Moonshine via fingerprint twining;
 \ClaimStatusConditional]
\label{thm:part-ii-moonshine-via-fingerprint}
Let $V^\natural$ be the Monster vertex operator algebra
\textup{(}Frenkel--Lepowsky--Meurman\textup{)}, realized as the
FLM $\mathbb{Z}/2$-orbifold extension
$V_{\Lambda_{\mathrm{Leech}}}^+\oplus
V_{\Lambda_{\mathrm{Leech}}}^{T,+}$ of the Leech lattice VOA\@.
The fingerprint $\varphi(V^\natural)$ is
\begin{equation}
 \varphi(V^\natural) \;=\;
 \bigl(\,4,\; r_{\max}^{\mathrm{res}}(V^\natural),\;
 \chi_{\VOA}(V^\natural),\;
 n_{\mathrm{Griess}}(V^\natural),\;
 \mathrm{coset}_{V^\natural}\,\bigr),
\end{equation}
recording $p_{\max}(V^\natural) = 4$ from the stress-tensor OPE,
$\kappa(V^\natural)=12$ from the central term, and a residue-depth
slot determined by $\mathsf H_{\mathrm{res}}(V^\natural;X)$, together
with
$n_{\mathrm{Griess}}(V^\natural)$ from the PBW-compatible
weight-$2$ Griess generators. The condition
$\dim V^\natural_1=0$ removes affine currents; it does not make
the Virasoro vector a strong generator of $V^\natural$.
The coset slot is
\begin{equation}
 \mathrm{coset}_{V^\natural} \;=\;
 \bigl(\,V_{\Lambda_{\mathrm{Leech}}},\; \mathbb{Z}/2,\;
 V^\natural_{\mathrm{tw}}\,\bigr)
\end{equation}
\textup{(}Leech lattice, orbifold involution, twisted sector\textup{)}.
For each $g \in \mathrm{Monster}$, the McKay--Thompson series
\begin{equation}
 T_g(\tau) \;=\; \mathrm{Tr}\bigl|_{V^\natural}\bigl(g\,q^{L_0 -
 1}\bigr)
\end{equation}
lifts to a twining fingerprint $\varphi_g(V^\natural)$, with
twining coset slot $\mathrm{OrbSec}_g(V^\natural)$, the
$g$-twisted orbifold-sector datum of the moonshine module. For
elements lying in the Leech-lattice subgroup this specializes to the
corresponding lattice fixed-sector notation; for a general Monster
element it is not a sub-VOA $V_{\Lambda_{\mathrm{Leech}}}^g$.
The assignment
$g \mapsto \varphi_g(V^\natural)$ is constant on Monster
conjugacy classes, yielding $194$ fingerprints indexed by
$\mathrm{Conj}(\mathrm{Monster})$; the frame shapes of these
$194$ twining fingerprints equal the Atkin--Lehner-twined bar
Euler characters
\textup{(}Theorem~\ref{thm:moonshine-bar-euler-master}\textup{)}.
\end{theorem}

\begin{proof}
Fingerprint of $V^\natural$: $p_{\max}(V^\natural) = 4$ by the
stress-tensor OPE
\textup{(}$V^\natural$ contains the Virasoro sub-VOA at $c = 24$,
whose $TT$ self-OPE has quartic pole\textup{)};
the residue-depth slot is computed from a compatible scalar
contraction on the inherited Virasoro sub-VOA and its embedding into
$V^\natural$; the local central coefficient gives
$\kappa(V^\natural)=c/2=12$;
$\chi_{\VOA}(V^\natural)$ is the Hodge-filtered chiral Euler
character, explicitly the $J$-function $J(\tau) = q^{-1} +
\textup{196884}\,q + \cdots$ \textup{(}normalized so that the
constant term vanishes\textup{)};
$n_{\mathrm{Griess}}(V^\natural)$ records the minimal
PBW-compatible strong generators in the weight-$2$ Griess algebra.
$V^\natural_1=0$ by Frenkel--Lepowsky--Meurman.  The Virasoro vector
generates its Virasoro sub-VOA, while PBW strong generation uses the
weight-two Griess algebra.  The fourth slot is therefore a
Griess-algebra invariant; and
$\mathrm{coset}_{V^\natural}$ encodes the Leech-lattice orbifold via
the triple $(V_\Lambda, \mathbb{Z}/2, V^\natural_{\mathrm{tw}})$.
The integer $196884$, the dimension of the weight-$2$ Griess
primary component modulo $\mathbb{C}\omega$ and equivalently the
coefficient of $q^1$ in $J(\tau)$, enters the fingerprint through
$\chi_{\VOA}$ and the Griess slot, not through a one-generator
Virasoro reduction.

Twining lift: for $g \in \mathrm{Monster}$, the $g$-twined bar
Euler character is
\begin{equation}
 \chi^g_{\VOA}(V^\natural) \;=\; \mathrm{Tr}\bigl|_{V^\natural}
 \bigl(g\,q^{L_0 - 1}\bigr) \;=\; T_g(\tau),
\end{equation}
by Theorem~\ref{thm:moonshine-bar-euler-master}
\textup{(}chiral-moonshine bar-Euler master identity\textup{)}; the
twining coset slot is the $g$-twisted orbifold-sector datum of
$V^\natural$. It reduces to a Leech-lattice fixed-sector notation
only for elements represented on the lattice carrier. The structural
slots $p_{\max}$, $r_{\max}$, and the Griess count are
$g$-independent; the Euler slot is the class function $T_g$.

Class-function property: the twining fingerprint depends on $g$
only through its Monster conjugacy class because $L_0$ commutes
with the Monster action, so $T_{h g h^{-1}}(\tau) = T_g(\tau)$ for
all $h \in \mathrm{Monster}$; the coset slot transforms
$\mathrm{Monster}$-equivariantly. There are
$|\mathrm{Conj}(\mathrm{Monster})| = 194$ classes; hence $194$
twining fingerprints.

Agreement with frame shapes: for each $g$, the frame shape of
$T_g(\tau)$ is the Atkin--Lehner-twined bar Euler character of
$V^\natural$
\textup{(}Theorem~\ref{thm:shadow-tower-twining-universality}\textup{)};
this is the fingerprint-level statement of the
Conway--Norton correspondence.
\end{proof}

The fingerprint twining assembles moonshine into a family of
$194$ twining fingerprints, one per Monster class,
recovered from the bar complex equivariantly.

\section{Higher-genus fingerprint}
\label{sec:part-ii-higher-genus}

The fingerprint extends coherently across the genus tower.

\begin{proposition}[Higher-genus fingerprint;
 \ClaimStatusConditional]
\label{prop:part-ii-higher-genus-fingerprint}
For each $g \ge 0$, the fingerprint $\varphi_g(\cA)$ at genus $g$
is obtained from $\varphi_0(\cA)$ by:
\begin{itemize}
 \item $p_{\max}$ unchanged at each genus \textup{(}pole order is
 an algebraic invariant, independent of the curve\textup{)};
 \item $r_{\max}$ unchanged at each genus \textup{(}shadow depth
 is genus-independent by Theorem~D and
 Corollary~\ref{cor:mumford-multiplicative}\textup{)};
 \item $\chi_{\VOA, g}(\cA) = \chi_{\VOA, 0}(\cA) \cdot
 (1 - q^g)^{-\kappa(\cA)}$ in the shadow-kernel
 \textup{(}modular-bootstrap closed form\textup{)};
 \item $n_{\mathrm{strong}}(\cA)$ genus-independent;
 \item $\mathrm{coset}_g(\cA)$ obtained from $\mathrm{coset}_0(\cA)$
 by clutching: the genus-$g$ coset is the $g$-fold sewing of
 the genus-$0$ coset along the sewing parameters
 $\{q_i\}_{i=1}^{3g-3}$ on $\overline{\mathcal{M}}_{g,n}$.
\end{itemize}
The genus tower of fingerprints $\{\varphi_g(\cA)\}_{g \ge 0}$ is
governed by the shadow obstruction tower
$\{S_r(\cA)\}_{r \ge 2}$ via the uniform-weight formula
the Hodge character
$\operatorname{ch}_g(\mathfrak O_g^K(\cA))=
(-1)^g\kappa(\cA)\lambda_g$ under~$H_D^K$ and the graph functional
$F_g=\kappa\lambda_g^{\mathrm{FP}}+\delta F_g^{\mathrm{cross}}$
under~$H_D^{\mathrm{graph}}$.
\end{proposition}

\begin{proof}
Pole order and shadow depth are algebraic invariants of $B(\cA)$,
independent of the curve on which chiral factorization homology
is computed; $\chi_{\VOA}$ and $\mathrm{coset}$ acquire genus
dependence through the perfect-object formula
$\mathfrak O_g^K(\cA)=\kappa(\cA)\lambda_{-1}(\mathbb E_g)$ under
$H_D^K$ and the graph correction
$\delta F_g^{\mathrm{cross}}$ under~$H_D^{\mathrm{graph}}$ of
Theorem~\ref{thm:multi-weight-genus-expansion}; the explicit
genus-$g$ modular-bootstrap formula for $\chi_{\VOA, g}$ is
computed in Chapter~\ref{chap:higher-genus}. The coset clutching
formula is Chapter~\ref{ch:genus-complete}
Proposition~\ref{prop:clutching-coset}.
\end{proof}

The higher-genus fingerprint is a corollary of the genus-$0$
fingerprint and the shadow tower; the same five slots carry the
higher-genus data.

\section{Typed endpoint: completeness}
\label{sec:part-ii-platonic-endpoint}

The five Part~II structural surfaces assemble into a single
conditional endpoint.

\begin{theorem}[Typed completeness of Part~II;
 \ClaimStatusConditional]
\label{thm:part-ii-platonic-completeness}
There is a canonical map
\begin{equation}
 \varphi\colon \mathrm{StdLand}^{\mathrm{chirKosz}}
 \;\longrightarrow\; \mathcal{F},
 \qquad \cA \;\longmapsto\;
 (p_{\max}, r_{\max}^{\mathrm{res}}, \chi_{\VOA}, n_{\mathrm{strong}},
 \mathrm{coset})
\end{equation}
from the standard landscape of chirally Koszul chiral algebras to
the fingerprint space $\mathcal{F}$ with the following properties.
\begin{enumerate}[(i)]
\item \textup{(}Conditional reconstruction;
 Theorem~\ref{thm:part-ii-fingerprint-reconstruction-surface}\textup{)}
 $\varphi$ is invariant under bar-complex isomorphism and becomes
 injective on bar-complex isomorphism classes on the residual-free
 fingerprint-reconstruction surface.
 \item \textup{(}Coarse projection;
 Theorem~\ref{thm:part-ii-quadrichotomy-coarse-projection}\textup{)}
 The four-class assignment
 $\Pi_{\mathrm{coarse}}: \mathcal{F} \to \{\mathrm{G}, \mathrm{L},
 \mathrm{C}, \mathrm{M}\}$ is a lossy quotient with explicit
 witnesses $H_k, V_k(\fsl_2), \beta\gamma, \Vir_c$. A fifth
 stratum $\mathrm{FF}$ appears canonically at $\kappa = 0$.
 \item \textup{(}Motivic comparison problem;
 Problem~\ref{thm:part-ii-motivic-purity-virasoro-r-le-8}\textup{)}
 The packages $\mathsf H_{\mathrm{res}}^{\le8}$ and
 $\mathsf H_{\mathrm{mot}}^{\le8}$ define the construction problem
 for motivic scalar residues through arity~$8$; Tate factorization and
 denominator divisors are named subproblems.
 \item \textup{(}Moonshine twining;
 Theorem~\ref{thm:part-ii-moonshine-via-fingerprint}\textup{)}
 The fingerprint $\varphi(V^\natural)$ of the Monster VOA lifts
 to $194$ twining fingerprints indexed by
 $\mathrm{Conj}(\mathrm{Monster})$, with frame shapes equal to
 bar Euler characters.
 \item \textup{(}Higher-genus coherence;
 Proposition~\ref{prop:part-ii-higher-genus-fingerprint}\textup{)}
 The genus tower $\{\varphi_g(\cA)\}_{g \ge 0}$ is governed by
 the shadow obstruction tower and UNIFORM/ALL-WEIGHT genus
 corrections.
\end{enumerate}
$\varphi$ is canonical \textup{(}naturally isomorphic for any
choice of coset representative\textup{)}, complete, coarsely
projective, motivically lifted, moonshine-equivariant, and
genus-coherent. It is the characteristic datum of Part~II.
\end{theorem}

\begin{proof}
Clauses (i)-(v) are the cited theorems; canonicity follows from
the fact that each slot is independent of the choices made in
its construction (pole order by minimal generator choice, shadow
tower by PBW filtration choice, chiral Euler character by
Hodge-filtration choice, strong-generator count by generator
minimization, coset by canonical commutant). Naturality under
chiral-algebra isomorphism is Proposition~\ref{prop:varphi-natural}.
\end{proof}

\section{The paramodular master over $\overline{\mathcal A_2}$}
\label{sec:part-ii-paramodular-master}

The one-coordinate reading of the characteristic datum reaches its
K3 apex at the master paramodular object. Write
$\mathfrak g_{\Delta_5} = \mathrm{Borch}(F_3,\phi^{K3}_{0,1})$ for
the K3-BKM Lie superalgebra produced by Borcherds lift from the
hyperbolic lattice $\Lambda_{\mathrm{Muk},II}^{2,1}$: a rank-$3$
hyperbolic core orthogonal to the rank-$24$ Mukai lattice. Its
compact Hall--Drinfeld recognition target
\(\mathbf H_{\Delta_5}^{\mathrm{tgt}}\) has Weyl denominator
\(\Delta_5 \in S_5(K(1))\), the Gritsenko paramodular weight-$5$
cusp form.  The notation \(\mathbf H_{\Delta_5}\) is reserved for the
recognized compact object carrying a finite Hall source, Hopf pairing,
source-to-target comparison maps, a PBW presentation theorem, and a
Mittag--Leffler inverse-limit passage.

\begin{theorem}[K3 determinant comparison problem for Part~II;
 \ClaimStatusOpen]
\label{thm:part-ii-k3-master-L-value}
Assume the compact Hall--Borcherds recognition package for the target
\(\mathbf H_{\Delta_5}^{\mathrm{tgt}}\).  The additional package
$H_{\Delta_5}^{\mathrm{det}}$ consists of the genus-$1$ determinant
line, its Quillen metric, a trace-class sewing operator, and a
comparison with the automorphic line generated by~$\Delta_5$.  The
required comparison is
\[
\log Z^{(1),\mathrm{sew}}_{\mathbf H_{\Delta_5}}
\xrightarrow{\ H_{\Delta_5}^{\mathrm{det}}\ }
-\log\bigl\|\Delta_5\bigr\|^2_{\mathrm{Quillen}}.
\]
The exact topological input is $\chi_{\mathrm{top}}(K3)=24$.
An $L$-value refinement requires a theorem identifying the regularized
determinant, the standard $L$-function normalization, and its
coefficient.  This is the open analytic obligation of
$H_{\Delta_5}^{\mathrm{det}}$; see
Chapter~\ref{V3-ch:k3-chiral-bialgebra-platonic}.
\end{theorem}

At the K3 base point the one-coordinate shadow $\kappa$ rises to
a five-coordinate fingerprint by inheriting the slots of
$\Psi_\bullet$, so the fingerprint theorem
(Theorem~\ref{thm:part-ii-fingerprint-reconstruction-surface}) reads
the Quillen norm slot-by-slot. The Baily--Borel--Freitag
stratification of $\overline{\mathcal A_2}$ produces four sibling
functors $\{\Psi,\Psi^{\deg},\Psi^{\mathrm{tor}},\Psi^{\mathrm{metap}}\}$
indexed by reflective interior, Klingen cusp (super-affine),
Humbert divisor (quantum-toroidal), metaplectic branch (Conway
$V^{s\natural}$). Under averaging these four paramodular siblings
project to the four shadow-tower archetypes $G/L/C/M$; the Humbert
divisor coincides with the Argyres--Douglas locus of the
class-$\mathcal S$ SCFT $\mathcal T[A_1,\Sigma_{0,24}]$, where the
Seiberg--Witten curve degenerates to $(A_1,A_{2N-1})$
(Gaiotto--Moore--Neitzke $2009$, Example~$8.3$). The fifth stratum
$FF$ of Theorem~\ref{thm:part-ii-quadrichotomy-coarse-projection}
is the critical-level degeneration of the reflective interior.

\medskip

\begin{remark}[Humbert--Heegner admissibility filtration on
the pentagon tower]
\label{rem:part-ii-humbert-heegner}
On the scalar Siegel--Borcherds target, any constructed pentagon tower
\(\{\phi^{(n)}\}_{n\ge 3}\) compatible with the Humbert--Heegner gates
must satisfy
\[
\phi^{(n)} \;=\; 0
\qquad\text{unless}\qquad
n \;\equiv\; 3,5 \pmod 8,
\qquad
D_n \;=\; (n-3)/2 \;\in\; \{0,1\}\pmod 4,
\]
and the first admissible non-vanishing rung is
$\phi^{(5)} = 2\,[\mathrm{gen}]^{\otimes 5}$ in the normalization of
Theorem~\ref{thm:bc-polar-depth-reduction}. For admissible
$n\ge 11$ the polar coefficient vanishes. The admissibility condition
on $(D_n\pmod 4)$ is a Heegner-divisor reading of the shadow
denominator pattern.  An all-orders associator on a Hall realisation
additionally requires the compact Hall source, comparison maps, and
pentagon identities on that source.
\end{remark}

\begin{remark}[Bloch--Kato Selmer deformation of the fingerprint]
\label{rem:part-ii-selmer-deformation}
The target fingerprint $\varphi(\mathbf H_{\Delta_5}^{\mathrm{tgt}})$ admits a
one-parameter arithmetic deformation: the Bloch--Kato Selmer group
$\dim H^1_f(\mathbb Q,\mathrm{std}\,\rho_{\Delta_5}) = 1$ supplies
the paramodular cyclotomic Hida-family tangent along which the
fingerprint coset slot varies. This one-dimensional Selmer line
is the arithmetic tangent space to the fingerprint map~$\varphi$
at the K3 base point.
\end{remark}

\begin{remark}[Metaplectic $\mathrm{grt}_1^{(1/2)}$ on the shadow tower]
\label{rem:part-ii-metaplectic-grt}
The K3 comparison asks for a morphism from the motivic residue package
of Theorem~\ref{thm:part-ii-motivic-purity-virasoro-r-le-8} to a
metaplectic motivic Lie algebra carrying the Saito--Kurokawa classes.
The target normalization may use the exact arithmetic identity
$\tau(2)^2/2=288$.  A bracket formula requires the comparison
morphism, its grading convention, and a chain representative for the
target classes.
\end{remark}

\section{Reading order within Part II}
\label{sec:part-ii-reading-order}

A reader approaching Part~II may follow two orders. The
\emph{fingerprint order} starts at
\S\ref{sec:part-ii-fingerprint} (conditional reconstruction) and descends
to the individual coordinates.  The \emph{residue order} starts with
the exact Virasoro OPE and Gram form, continues with
$\mathsf H_{\mathrm{res}}$, and then introduces the motivic and
asymptotic comparison packages.

We recommend the residue order for the first reading: the central OPE
coefficient and the level-four Gram form are concrete, while
Chapter~\ref{ch:shadow-tower-higher-coefficients} states the precise
chain construction for higher scalar coordinates.  The fingerprint
order is
suitable for a second reading, when the fingerprint $\mathcal{F}$
picture guides the interpretation of every structural invariant as
a fingerprint slot.

\medskip

The characteristic datum of Part~II is the fingerprint.  The central
OPE coefficient supplies its local scalar input; the residue package
constructs higher scalar coordinates; the motivic and denominator
packages formulate the arithmetic comparison.

\section{Dependency map}
\label{sec:part-ii-dependency-map}

The five structural surfaces of Part~II rest on concrete local
anchors and named comparison packages.

\begin{center}
\begin{tabular}{p{3.5cm}|p{6cm}|p{5cm}}
\hline
\textbf{Structural theorem surface} & \textbf{Anchor chapters (Part II)}
 & \textbf{Cross-volume anchors} \\
\hline
Theorem~\ref{thm:part-ii-fingerprint-reconstruction-surface}
(conditional fingerprint reconstruction)
 & Chapter~\ref{ch:infinite-fingerprint}
 & Vol II
 Infinite Fingerprint; Vol III
 $\varphi(K3)$, $\varphi(X_{\mathrm{CY}_3})$. \\
\hline
Theorem~\ref{thm:part-ii-quadrichotomy-coarse-projection}
(coarse projection)
 & Chapter~\ref{ch:infinite-fingerprint}
 \S\ref{sec:coarse-projection-fibres}
 & Vol II Infinite Fingerprint; Vol III Calabi--Yau fingerprint
 stratification. \\
\hline
Problem~\ref{thm:part-ii-motivic-purity-virasoro-r-le-8}
(motivic residue problem $r \le 8$)
 & Chapter~\ref{ch:virasoro-motivic-purity};
 Chapter~\ref{ch:shadow-tower-higher-coefficients}
 & Exact Ward correlators and level-four Gram form;
 $\mathsf H_{\mathrm{res}}^{\le8}$,
 $\mathsf H_{\mathrm{mot}}^{\le8}$, and
 $\mathsf H_{\mathrm{Tate}}^{\le8}$ are the named construction
 packages. \\
\hline
Theorem~\ref{thm:part-ii-moonshine-via-fingerprint}
(moonshine twining)
 & Chapter~\ref{chap:chiral-moonshine-unified}
 & Borcherds denominator;
 FLM Monster VOA construction; Conway--Norton; the Vol.~II
 \(E_3^{\mathrm{top}}\) comparison for \(V^\natural\). \\
\hline
Theorem~\ref{thm:part-ii-platonic-completeness}
(structural completeness)
 & assembled from (i)-(v)
 & all. \\
\hline
\end{tabular}
\end{center}

\section{Claim status on the five structural surfaces}
\label{sec:part-ii-hz-iv-closures}

The five structural surfaces of Part~II carry one mathematical status each.
Proved rows name their derivations; open rows name the missing construction
package.

\begin{center}
\begin{tabular}{p{4cm}|p{5cm}|p{6cm}}
\hline
\textbf{Surface} & \textbf{Input} & \textbf{Evidence or remaining
obligation} \\
\hline
Theorem~\ref{thm:part-ii-fingerprint-reconstruction-surface}
(conditional fingerprint reconstruction)
 & Bar-complex isomorphism invariants
 \textup{(}pole order, shadow depth, $\chi_{\VOA}$,
 $n_{\mathrm{strong}}$, $\mathrm{coset}$\textup{)}
 & Slot independence is
 Theorem~\ref{thm:pole-depth-independence} and
 Proposition~\ref{prop:chi-strong-coset-independence}; conditional
 reconstruction is
 Theorem~\ref{thm:fingerprint-reconstruction-equivalence}. \\
\hline
Theorem~\ref{thm:part-ii-quadrichotomy-coarse-projection}
(coarse projection)
 & Shadow tower
 depth-indexed classification
 \textup{(}Chapter~\ref{ch:infinite-fingerprint}\textup{)}
 & Creutzig--Linshaw class
 assignment for $\cW$-algebra families; independent
 landscape census. \\
\hline
Problem~\ref{thm:part-ii-motivic-purity-virasoro-r-le-8}
(motivic residue construction)
 & Ordered Ward correlator complex and
 $\mathsf H_{\mathrm{res}}^{\le8}$
 & Motivic period comparison, ordered-to-symmetric averaging, and
 Tate-factorization checks; exact local anchor
 $\det G^{\Vir}_{4}=c^2(5c+22)/2$. \\
\hline
Theorem~\ref{thm:part-ii-moonshine-via-fingerprint}
(moonshine twining)
 & Frenkel--Lepowsky--Meurman
 $V^\natural$ construction; McKay--Thompson
 series; bar Euler master identity
 & Borcherds
 denominator identity route; Conway--Norton verification
 against $194$ Monster classes; independent ATLAS conjugacy
 class lookup. \\
\hline
Theorem~\ref{thm:part-ii-platonic-completeness}
(completeness)
 & Structural surfaces (i)--(v) assembled
 & Logical conjunction of the five preceding structural statements,
 with their stated hypothesis packages. \\
\hline
\end{tabular}
\end{center}

The general statements use the cited theorem proofs and the named open
packages; this table adds no independent reconstruction claim.

\medskip

The structural form of Part~II is specified. The existing chapters
of the Characteristic Datum are the concrete realizations of the
five slots; this introduction orients the reader and the subsequent
arguments.
